Let $n\ge 2$, $1\le p\le\infty$, and let $f\in L^p(\mathbb R^n)$ be arbitrary. For $r>0$ let $B_r^\alpha$ be the Fourier multiplier on $\mathbb R^n$ with symbol $(1-|\xi|^2/r^2)_+^\alpha$, and put
\[
\delta(p)=\max\{0,n|1/p-1/2|-1/2\}.
\]
If $\alpha>\delta(p)$, then
\[
\sup_{r>0}\|B_r^\alpha f\|_{L^p(\mathbb R^n)}
\le C_{n,p,\alpha}\|f\|_{L^p(\mathbb R^n)}.
\]
Under the same condition on $\alpha$, the means converge to $f$ in $L^p$ as $r\to\infty$ if $p<\infty$. If $2\le p<\infty$ and $\alpha>\max\{0,(n-1)/2-n/p\}$, then also
\[
\left\|\sup_{r>0}|B_r^\alpha f|\right\|_{L^p(\mathbb R^n)}
\le C_{n,p,\alpha}\|f\|_{L^p(\mathbb R^n)},
\]
and the means converge to $f$ almost everywhere as $r\to\infty$. The supremum is taken for a representative continuous in $r>0$ almost everywhere.
